\enquestion{8.4}{
    A confidence interval estimate is desired for the gain in a circuit on a semiconductor device.
    Assume that gain is normally distributed with standard deviation $\sigma = 25$.
    \begin{enumerate}[label=(\alph*)]
        \item Find a $\qty{95}{\percent}$ CI for $\mu$ when $n = 10$ and $\bar{x} = 1000$.
        \item Find a $\qty{95}{\percent}$ CI for $\mu$ when $n = 25$ and $\bar{x} = 1000$.
        \item Find a $\qty{99}{\percent}$ CI for $\mu$ when $n = 10$ and $\bar{x} = 1000$.
        \item Find a $\qty{99}{\percent}$ CI for $\mu$ when $n = 25$ and $\bar{x} = 1000$.
        \item How does the length of the CIs computed above change with the changes in sample size and confidence level?
    \end{enumerate}
}

\enquestion{8.7}{
    Consider the gain estimation problem in Exercise 8-4.
    \begin{enumerate}[label=(\alph*)]
        \item How large must $n$ be if the length of the $\qty{95}{\percent}$ CI is to be $30$?
        \item How large must $n$ be if the length of the $\qty{99}{\percent}$ CI is to be $30$?
    \end{enumerate}
}

\enquestion{8.9}{
    Suppose that $n = 100$ random samples of water from a freshwater lake were taken and the calcium concentration ($\qty{}{\mg\per\L}$) measured.
    A $\qty{95}{\percent}$ CI on the mean calcium concentration is $0.52 \le \mu \le 0.74$.
    \begin{enumerate}[label=(\alph*)]
        \item Would a $\qty{99}{\percent}$ CI calculated from the same sample data be longer or shorter?
        \item Consider the following statement: \enquote{There is a $\qty{95}{\percent}$ chance that $\mu$ is between $0.52$ and $0.74$.} Is this statement correct?
        Explain your answer.
        \item Consider the following statement: \enquote{If $n = 100$ random samples of water from the lake were taken and the $\qty{95}{\percent}$ CI on $\mu$ computed, and this process were repeated $1000$ times, $950$ of the CIs would contain the true value of $\mu$.} Is this statement correct?
        Explain your answer.
    \end{enumerate}
}

\enquestion{8.11}{
    The yield of a chemical process is being studied.
    From previous experience, yield is known to be normally distributed and $\sigma = 3$.
    The past five days of plant operation have resulted in the following percent yields: $91.6$, $87.34$, $90.8$, $89.65$, and $91.3$.
    Find a $\qty{95}{\percent}$ two-sided confidence interval on the true mean yield.
}

\enquestion{8.17}{
    The life in hours of a $75$-watt light bulb is known to be normally distributed with $\sigma=20$ hours.
    % A random sample of $n$ bulbs has a mean life of $\bar{x}=1014$ hours.
    Suppose you wanted the total width of the two-sided confidence interval on mean life to be six hours at \qty{95}{\percent} confidence.
    What sample size should be used?
}

\enquestion{8.29}{
    A postmix beverage machine is adjusted to release a certain amount of syrup into a chamber where it is mixed with carbonated water.
    A random sample of $25$ beverages was found to have a mean syrup content of $\bar{x} = \qty{30}{\mL}$.
    Suppose the standard deviation is known and equal to $\sigma = \qty{0.5}{\mL}$.
    Find a $\qty{95}{\percent}$ CI on the mean volume of syrup dispensed.
}

\enquestion{7.63}{
    Let $f(x) = (1 / \theta) x^{(1 - \theta) / \theta}$, $0 < x < 1$, and $0 < \theta < \infty$.
    Show that $\hat{\Theta} = -(1 / n) \sum_{i=1}^{n} \ln(X_{i})$ is the maximum likelihood estimator for $\theta$ and that $\hat{\Theta}$ is an unbiased estimator for $\theta$.
}